% ============================================================
%  preamble.tex  --  packages, environments, boxes
% ============================================================
\usepackage[utf8]{inputenc}
\usepackage[T1]{fontenc}
\usepackage[expansion=false]{microtype}
\usepackage{amsmath,amssymb,amsthm,mathtools,bm}
\usepackage{mathrsfs}
\usepackage[sc]{mathpazo}
\usepackage{courier}
\usepackage[paperwidth=215mm,paperheight=285mm,inner=21.8mm,textwidth=111.4mm,marginparsep=8.2mm,marginparwidth=50.5mm,top=30mm,textheight=208mm,headheight=12pt,headsep=14mm]{geometry}
\usepackage{sidenotes}
\usepackage[font=small,labelfont=bf,labelsep=period]{caption}
\usepackage{enumitem}
\usepackage{graphicx}
\usepackage{makeidx}
\usepackage{longtable}
\makeindex
\emergencystretch=4em
\usepackage{booktabs}
\usepackage{xcolor}
\usepackage{tikz}
\usepackage{tikz-cd}
\usetikzlibrary{arrows.meta,positioning,calc,decorations.pathreplacing}
\usepackage[most]{tcolorbox}
\usepackage{fancyhdr}
\usepackage{titlesec}
\usepackage[numbers,sort&compress]{natbib}
\usepackage[colorlinks=true,linkcolor=physcol!75!black,citecolor=physcol!75!black,urlcolor=physcol!75!black]{hyperref}
\usepackage[nameinlink,capitalise]{cleveref}

\setlength{\parskip}{0pt}
\setlength{\parindent}{1.45em}
\linespread{1.05}
\setcounter{tocdepth}{1}
\setcounter{secnumdepth}{2}
\renewcommand{\footnote}[1]{\sidenote{#1}}
\newcommand{\oldnum}[1]{{\oldstylenums{#1}}}

% ---------- colours ----------
\definecolor{physcol}{RGB}{30,90,160}
\definecolor{warncol}{RGB}{170,40,40}
\definecolor{dictcol}{RGB}{40,120,70}
\definecolor{statuscol}{RGB}{140,90,0}

% ---------- theorem environments ----------
\theoremstyle{plain}
\newtheorem{theorem}{Theorem}[chapter]
\newtheorem{proposition}[theorem]{Proposition}
\newtheorem{lemma}[theorem]{Lemma}
\newtheorem{corollary}[theorem]{Corollary}
\theoremstyle{definition}
\newtheorem{definition}[theorem]{Definition}
\newtheorem{example}[theorem]{Example}
\newtheorem{exercise}{Exercise}[section]
\theoremstyle{remark}
\newtheorem{remark}[theorem]{Remark}

\crefname{theorem}{Theorem}{Theorems}
\crefname{exercise}{Exercise}{Exercises}

% proof sketch (the book mostly sketches)
\newenvironment{proofsketch}{\begin{proof}[Idea of proof]}{\end{proof}}

% ---------- boxes: quiet run-in notes with a thin coloured rule ----------
\tcbset{notebox/.style={enhanced,breakable,frame hidden,boxrule=0pt,sharp corners,
  left=9pt,right=4pt,top=5pt,bottom=5pt,before skip=9pt,after skip=9pt,
  attach title to upper={\ \ },fonttitle=\scshape,coltitle=black,
  colbacktitle=white,before upper={\parindent=0pt\parskip=3pt\relax}}}
\usepackage{environ}
% the physics motivation lives in the outer margin, like a sidenote
\NewEnviron{physmotiv}[1][Physics motivation]{%
  \marginpar{\footnotesize\raggedright\setlength{\parindent}{0pt}\setlength{\parskip}{3pt}\linespread{1.0}\selectfont
  \textsc{#1.}\ \BODY}}
\newtcolorbox{whatwrong}[1][What could go wrong]{notebox,colback=warncol!6,
  borderline west={1.6pt}{0pt}{warncol},title={#1.}}
\newtcolorbox{dictbox}[1][Physics $\leftrightarrow$ Math dictionary]{notebox,colback=dictcol!6,
  borderline west={1.6pt}{0pt}{dictcol},title={#1.}}
\newtcolorbox{unsettled}[1][Status of the literature]{notebox,colback=statuscol!7,
  borderline west={1.6pt}{0pt}{statuscol},title={#1.}}
\newtcolorbox{keyidea}[1][Key idea]{notebox,colback=black!5,
  borderline west={1.6pt}{0pt}{black!65},title={#1.}}

% Section marker for skippable material
\newcommand{\skippable}{\par\noindent{\small\itshape (This section can be skipped on a first reading; see the reading paths in \cref{sec:paths}.)}\par}

% ---------- running heads and headings (after Schwichtenberg, Physics from Symmetry) ----------
\newcommand{\booktitleshort}{operator algebras and quantum gravity}
\pagestyle{fancy}
\fancyhf{}
\renewcommand{\chaptermark}[1]{\markboth{#1}{#1}}
\renewcommand{\sectionmark}[1]{}
\fancyhead[LE]{\small\textsc{\thepage\hspace{1.2em}\booktitleshort}}
\fancyhead[RO]{\small\textsc{\nouppercase{\rightmark}\hspace{1.2em}\thepage}}
\renewcommand{\headrulewidth}{0pt}
\fancypagestyle{plain}{\fancyhf{}\renewcommand{\headrulewidth}{0pt}}

\renewcommand{\thepart}{\ifnum\value{part}=0 0\else\Roman{part}\fi}
\titleformat{\part}[display]{\normalfont\bfseries}{\fontsize{22}{26}\selectfont\partname\ \oldnum{\thepart}}{14pt}{\fontsize{30}{34}\selectfont}
\titleformat{\chapter}[display]{\normalfont\bfseries}{\fontsize{24}{28}\selectfont\oldnum{\thechapter}}{12pt}{\fontsize{27}{31}\selectfont}
\titleformat{name=\chapter,numberless}[display]{\normalfont\bfseries}{}{0pt}{\fontsize{27}{31}\selectfont}
\titlespacing*{\chapter}{0pt}{55pt}{34pt}
\titleformat{\section}{\normalfont\fontsize{14}{17}\selectfont\bfseries}{\oldnum{\thesection}}{0.9em}{}
\titlespacing*{\section}{0pt}{22pt plus 4pt minus 2pt}{8pt plus 2pt}
\titleformat{\subsection}{\normalfont\fontsize{11.5}{14}\selectfont\bfseries}{\oldnum{\thesubsection}}{0.8em}{}
\titlespacing*{\subsection}{0pt}{15pt plus 3pt minus 2pt}{5pt plus 1pt}
\titleformat{\subsubsection}[runin]{\normalfont\bfseries}{}{0pt}{}[.\ ]
\titlespacing*{\subsubsection}{0pt}{9pt plus 2pt}{0pt}
\captionsetup{figurewithin=chapter}
